% © 2026 Ing. David Jaroš — CC BY-NC-ND 4.0
%
% This work is licensed under a Creative Commons Attribution-NonCommercial-NoDerivatives
% 4.0 International License (CC BY-NC-ND 4.0).
%
% License History: Earlier drafts (up to v0.3) were released under CC BY 4.0.
% From v0.4 onward, all material is released under CC BY-NC-ND 4.0 to protect
% the integrity of the theoretical work during ongoing academic development.
%
% See LICENSE.md for full license text.

\ifdefined\INCLUDEMODE
\else
  \documentclass[12pt,a4paper]{article}
  \usepackage[a4paper, margin=2.5cm]{geometry}
  \usepackage{amsmath, amssymb, amsthm}
  \usepackage{mathtools}
  \usepackage{hyperref}
  \usepackage{xcolor}
  \usepackage{mdframed}
  \usepackage{booktabs}
  \usepackage{longtable}
  \usepackage{tcolorbox}

  \newtheorem{theorem}{Theorem}[section]
  \newtheorem{lemma}[theorem]{Lemma}
  \newtheorem{proposition}[theorem]{Proposition}
  \newtheorem{corollary}[theorem]{Corollary}
  \theoremstyle{definition}
  \newtheorem{definition}[theorem]{Definition}
  \theoremstyle{remark}
  \newtheorem{remark}[theorem]{Remark}

  \newmdenv[backgroundcolor=yellow!20, linecolor=orange, linewidth=1.5pt]{gapbox}
  \newmdenv[backgroundcolor=green!15, linecolor=green!60!black, linewidth=1.5pt]{resultbox}
  \newmdenv[backgroundcolor=red!8, linecolor=red!60, linewidth=1.5pt]{warnbox}
  \newmdenv[backgroundcolor=blue!8, linecolor=blue!50, linewidth=1.5pt]{insightbox}

  \title{\textbf{Symmetry-Breaking Projection Route to $\alpha$}\\
         \large UBT Alpha Derivation — Route A2}
  \author{Ing.\ David Jaro\v{s}}
  \date{2026-04-27}

  \begin{document}
  \maketitle
\fi

%──────────────────────────────────────────────────────────────────────────────
\section{Objective and Scope}
\label{sec:sb:objective}
%──────────────────────────────────────────────────────────────────────────────

This document attempts to derive the fine-structure constant~$\alpha$ from the
spontaneous symmetry breaking (SSB) of $SU(2)_L \times U(1)_Y \to U(1)_{\mathrm{EM}}$
within UBT.

The central question is whether the UBT biquaternion field $\Theta$ has a vacuum
structure that fixes the Weinberg angle $\theta_W$, thereby determining the
electromagnetic coupling $e = g\sin\theta_W$ without free parameters.

\begin{tcolorbox}[colback=blue!4!white, colframe=blue!50!black, title=Route A2 Status]
\begin{tabular}{ll}
Overall status & \textbf{CONDITIONAL} \\
Blocking gap & Gap EW-1 ($g'/g$ ratio not derived from $\mathcal{B}$) \\
Gap EW-2 & UBT derivation of the doublet VEV structure \\
What is proved & Standard electroweak algebra (not UBT-specific) \\
What remains & Fix $\theta_W$ from $\mathrm{Aut}(\mathcal{B})$ \\
\end{tabular}
\end{tcolorbox}

%──────────────────────────────────────────────────────────────────────────────
\section{Explicit Assumptions}
\label{sec:sb:assumptions}
%──────────────────────────────────────────────────────────────────────────────

\begin{enumerate}
  \item[\textbf{A1}] The UBT action contains $SU(2)_L \times U(1)_Y$ as a gauge
    symmetry at high energy.  \textit{Status}: canonical input.

  \item[\textbf{A2}] The Θ-field admits a non-zero vacuum expectation value
    $\Theta_0 \neq 0$ that is invariant under $U(1)_{\mathrm{EM}}$ but not under
    $SU(2)_L \times U(1)_Y$.  \textit{Status}: \textbf{assumed} (Gap EW-2).

  \item[\textbf{A3}] The VEV $\Theta_0$ transforms as a doublet of $SU(2)_L$ with
    hypercharge $Y = 1/2$, analogously to the SM Higgs doublet.
    \textit{Status}: \textbf{assumed} (not yet derived from $S[\Theta]$).

  \item[\textbf{A4}] The SSB pattern is the standard electroweak one:
    $SU(2)_L \times U(1)_Y \to U(1)_{\mathrm{EM}}$, with three eaten Goldstone modes
    giving mass to $W^\pm$, $Z^0$.  \textit{Status}: \textbf{assumed}.

  \item[\textbf{A5}] Natural units $\hbar = c = 1$.
\end{enumerate}

%──────────────────────────────────────────────────────────────────────────────
\section{Electroweak Symmetry Breaking in UBT}
\label{sec:sb:ewsb}
%──────────────────────────────────────────────────────────────────────────────

\subsection{The Higgs Analogue in UBT}

In the Standard Model, the Higgs doublet $H$ acquires a VEV
$\langle H \rangle = (0,\, v/\sqrt{2})^T$, breaking $SU(2)_L \times U(1)_Y$.

In UBT, the candidate for the Higgs role is the lowest $\psi$-mode of $\Theta$:
\begin{equation}
  \Theta_0(q) := \Theta(q, \tau)\big|_{\psi=0}
  \;\in\; \mathbb{C}\otimes\mathbb{H}.
  \label{eq:sb:Theta0}
\end{equation}
If $\Theta_0$ is required to be a singlet under $U(1)_{\mathrm{EM}}$ but a doublet
under $SU(2)_L$, the SSB occurs when $\Theta_0 \neq 0$.

\begin{gapbox}
\textbf{Gap EW-2}: Derive from $S[\Theta]$ that the potential $V(\Theta)$ produces
a non-zero $\Theta_0$ with doublet quantum numbers.  The Mexican-hat structure
$V(\Theta) = -\mu^2 \mathrm{Tr}[\Theta^\dagger\Theta] + \lambda(\mathrm{Tr}[\Theta^\dagger\Theta])^2$
gives a VEV $|\Theta_0|^2 = \mu^2/(2\lambda)$, but the representation of $\Theta_0$
under $SU(2)_L$ must be derived, not assumed.
\end{gapbox}

\subsection{Unbroken Generator}

The unbroken $U(1)_{\mathrm{EM}}$ generator is:
\begin{equation}
  Q = T_3 + \frac{Y}{2},
  \label{eq:sb:Q}
\end{equation}
satisfying $Q\,\Theta_0 = 0$ (the VEV is electrically neutral).

For the doublet $\Theta_0 = (0,\, v/\sqrt{2})^T$:
\begin{equation}
  T_3\,\Theta_0 = -\frac{1}{2}\Theta_0,
  \quad
  \frac{Y}{2}\,\Theta_0 = +\frac{1}{2}\Theta_0,
  \quad
  Q\,\Theta_0 = 0.
  \label{eq:sb:neutrality}
\end{equation}

This confirms that the lower component of the doublet is electrically neutral,
consistent with $U(1)_{\mathrm{EM}}$ being unbroken.

\subsection{Physical Gauge Bosons after SSB}

The three massless Goldstone modes $\pi^{\pm}$, $\pi^0$ are absorbed into
$W^\pm$, $Z^0$.  The remaining massless gauge boson is the photon:
\begin{align}
  W^\pm_\mu &= \frac{1}{\sqrt{2}}(W^1_\mu \mp i W^2_\mu),
  \label{eq:sb:Wpm} \\
  A_\mu     &= \sin\theta_W\, W^3_\mu + \cos\theta_W\, B_\mu,
  \label{eq:sb:Agamma} \\
  Z_\mu     &= \cos\theta_W\, W^3_\mu - \sin\theta_W\, B_\mu,
  \label{eq:sb:Z}
\end{align}
with the Weinberg angle determined by the coupling ratio:
\begin{equation}
  \tan\theta_W = \frac{g'}{g}.
  \label{eq:sb:tanW}
\end{equation}

%──────────────────────────────────────────────────────────────────────────────
\section{The Key Question: Does UBT Fix $\theta_W$?}
\label{sec:sb:fixing}
%──────────────────────────────────────────────────────────────────────────────

Equation~\eqref{eq:sb:tanW} shows that $\theta_W$ is fixed if and only if
the ratio $g'/g$ is fixed.  This is the crux of Route A2.

\subsection{Algebra Structure of $\mathcal{B} = \mathbb{C}\otimes\mathbb{H}$}

The biquaternion algebra has the following relevant substructure:
\begin{equation}
  \mathcal{B} \;\cong\; \mathbb{C} \otimes \mathbb{H}
  \;\cong\; M_2(\mathbb{C}),
  \label{eq:sb:B_iso}
\end{equation}
where $M_2(\mathbb{C})$ is the algebra of $2\times 2$ complex matrices.
Under this isomorphism:
\begin{itemize}
  \item The quaternion imaginary units $\mathbf{i}, \mathbf{j}, \mathbf{k}$ map to
        $i\sigma_1, i\sigma_2, i\sigma_3$ (with $\sigma_i$ Pauli matrices).
  \item The $\mathfrak{su}(2)_L$ generators are $\tau^i = \sigma^i/2$.
  \item The $\mathfrak{u}(1)_Y$ generator $Y$ must be an endomorphism of $M_2(\mathbb{C})$
        commuting with the $\mathfrak{su}(2)_L$ action.
\end{itemize}

By Schur's lemma, the only endomorphisms of $M_2(\mathbb{C})$ commuting with
$\mathfrak{su}(2)_L$ are multiples of the identity.  Therefore $Y \propto \mathbf{1}$
on each irreducible $SU(2)_L$ representation.

\begin{insightbox}
\textbf{Structural observation}: The hypercharge $Y$ is proportional to the identity
on each $SU(2)_L$ representation.  This is consistent with the SM assignment
$Y = 1/2$ for the Higgs doublet.  It does \emph{not} fix the magnitude of $g'$
relative to $g$, since the overall normalization of $Y$ is a free parameter.
\end{insightbox}

\begin{gapbox}
\textbf{Gap EW-1}: In $\mathcal{B} = M_2(\mathbb{C})$, the ratio of the $U(1)_Y$
generator norm to the $SU(2)_L$ generator norm is:
\begin{equation}
  \frac{\|Y\|}{\|\tau^3\|} = \frac{|y|}{\tfrac{1}{2}}  = 2|y|,
  \label{eq:sb:normratio}
\end{equation}
where $y$ is the hypercharge quantum number.  This ratio depends on the
hypercharge assignment, which is a free parameter in the UBT Lagrangian as
currently formulated.

\medskip
To fix $\theta_W$, one needs a principle from UBT that determines $g'$ relative to $g$.
No such principle has been identified from $S[\Theta]$ in the literature.
\end{gapbox}

\subsection{Attempt: Norm-Matching Condition}

One possible approach: demand that the kinetic term produces a unit-normalized
gauge kinetic term for both $W^i_\mu$ and $B_\mu$:
\begin{equation}
  -\frac{1}{4g^2} W^i_{\mu\nu} W^{i\,\mu\nu}
  - \frac{1}{4{g'}^2} B_{\mu\nu} B^{\mu\nu}
  \;\longrightarrow\;
  -\frac{1}{4}F^a_{\mu\nu} F^{a\,\mu\nu},
\end{equation}
with $F^a_{\mu\nu}$ in the canonically normalized frame.

If UBT requires the kinetic terms to be normalized identically for all gauge
bosons in a grand-unified sense, one would demand $g = g'$.  This gives:
\begin{equation}
  \theta_W\big|_{g=g'} = 45°,
  \qquad
  \sin^2\theta_W = \frac{1}{2},
  \qquad
  e = \frac{g}{\sqrt{2}}.
  \label{eq:sb:unification}
\end{equation}

\begin{warnbox}
\textbf{This prediction is incorrect.}
The observed value is $\sin^2\theta_W \approx 0.231 \neq 0.5$ at the $Z$-pole.
Therefore, $g = g'$ is not the right UBT constraint.  The norm-matching condition
does not work without renormalization group running.

If $g = g'$ at some UV unification scale $\Lambda_{\mathrm{GUT}}$, then running
down to $m_Z$ using SM beta functions gives $\sin^2\theta_W(m_Z) \approx 0.231$.
But this requires the SM running as additional input, and $\Lambda_{\mathrm{GUT}}$
is a new free parameter.  \emph{This is not a zero-parameter UBT derivation.}
\end{warnbox}

\subsection{Grand Unification Route (Possible Future Direction)}

If UBT embeds $SU(3)_c \times SU(2)_L \times U(1)_Y$ into a larger simple group
$G_{\mathrm{GUT}}$ (e.g., $SU(5)$, $SO(10)$, or an exceptional group) at some scale,
then the ratio $g'/g$ is fixed at that scale by the Lie algebra structure.

In $SU(5)$: $\sin^2\theta_W^{(0)} = 3/8$ at the GUT scale (algebraic result), which
runs to $\approx 0.231$ at $m_Z$.

\begin{insightbox}
\textbf{Possible Route A2 completion}: If UBT forces $G_{\mathrm{GUT}} = SU(5)$ or
a specific group that contains $SU(3) \times SU(2) \times U(1)$ as a maximal
subgroup, the relation $\sin^2\theta_W^{(0)} = 3/8$ would be a UBT prediction,
not a fitted input.  Whether $\mathcal{B}$ enforces this is an open algebraic question.
\end{insightbox}

%──────────────────────────────────────────────────────────────────────────────
\section{Partial Result: What Is Fixed by Symmetry Breaking}
\label{sec:sb:partial}
%──────────────────────────────────────────────────────────────────────────────

\begin{resultbox}
\textbf{What Route A2 establishes (conditionally):}
\begin{enumerate}
  \item The unbroken generator $Q = T_3 + Y/2$ and the photon field~\eqref{eq:sb:Agamma}
        are structurally correct given the standard SSB pattern.
  \item The norm-matching condition $g = g'$ is \emph{excluded} as a UBT constraint
        at the electroweak scale.
  \item A grand-unification completion of UBT within a simple group $G_{\mathrm{GUT}}$
        would fix $\sin^2\theta_W^{(0)}$ algebraically and, after RGE running,
        predict $\alpha$ at any scale.
  \item The VEV structure (Gap EW-2) is a necessary prior step.
\end{enumerate}
\end{resultbox}

%──────────────────────────────────────────────────────────────────────────────
\section{Failure Modes and Near-Misses}
\label{sec:sb:failures}
%──────────────────────────────────────────────────────────────────────────────

\begin{longtable}{p{3.5cm} p{7cm} p{3cm}}
\toprule
\textbf{Failure mode} & \textbf{Description} & \textbf{Assessment} \\
\midrule
$g = g'$ at EW scale &
  Predicts $\sin^2\theta_W = 1/2$, excluded by experiment. &
  Excluded \\
\midrule
No doublet structure &
  $\Theta_0$ may not transform as a doublet; SSB pattern may differ from SM. &
  Open question (Gap EW-2) \\
\midrule
No grand-unification &
  $\mathcal{B}$ may not embed into a simple $G_{\mathrm{GUT}}$; the ratio $g'/g$
  would then be genuinely free. &
  Most likely outcome \\
\midrule
Scale dependence &
  Even if $\theta_W$ is fixed at one scale, $\alpha$ at another scale requires
  RGE running, introducing additional inputs ($m_e$, $m_Z$). &
  Known limitation \\
\bottomrule
\end{longtable}

%──────────────────────────────────────────────────────────────────────────────
\section{Summary}
\label{sec:sb:summary}
%──────────────────────────────────────────────────────────────────────────────

\begin{center}
\begin{tabular}{llll}
\toprule
Result & Status & Gap & Source \\
\midrule
$Q = T_3 + Y/2$ (unbroken generator) & Proved (algebra) & — & \S\ref{sec:sb:ewsb} \\
Photon field~\eqref{eq:sb:Agamma} & Proved (algebra) & — & \S\ref{sec:sb:ewsb} \\
$\Theta_0$ as doublet VEV & Open & EW-2 & \S\ref{sec:sb:ewsb} \\
$g = g'$ at EW scale & \textbf{Excluded} & — & \S\ref{sec:sb:fixing} \\
$\tan\theta_W = g'/g$ from $\mathcal{B}$ & Open & EW-1 & \S\ref{sec:sb:fixing} \\
$G_{\mathrm{GUT}}$ completion & Possible future direction & — & \S\ref{sec:sb:fixing} \\
$\alpha$ from Route A2 & Conditional & EW-1, EW-2 & \S\ref{sec:sb:alpha} \\
\bottomrule
\end{tabular}
\end{center}

\noindent\textbf{Classification}: \textbf{CONDITIONAL} — blocked by Gaps EW-1 and EW-2.

\ifdefined\INCLUDEMODE
\else
  \end{document}
\fi
